\documentclass[a4paper,11pt]{article}
\usepackage[table]{xcolor}
\usepackage{amsfonts}
\usepackage{amssymb}
\usepackage{amsmath}
\usepackage{euscript}
\usepackage{boxedminipage}
\usepackage{lscape}
\usepackage{minitoc}
\usepackage{epsfig,psfrag,graphicx,verbatim}
\usepackage{natbib}
\usepackage{booktabs, longtable}
\usepackage{url}
\usepackage{makeidx}

%% Define a new 'code' style for the url package that will use a sans serif font.
\makeatletter
\def\url@codestyle{%
  \@ifundefined{selectfont}{\def\UrlFont{\sf}}{\def\UrlFont{\sffamily}}}
\makeatother
%% Now actually use the newly defined style.
\urlstyle{code}

\vfuzz2pt % Don't report over-full v-boxes if over-edge is small
\hfuzz2pt % Don't report over-full h-boxes if over-edge is small
\setlength{\oddsidemargin}{0cm}
\setlength{\evensidemargin}{0cm}
\setlength{\textwidth}{16cm}
\setlength{\textheight}{23cm}
\title{Test Model - Xparser 0.17.1 Circle (x,y)-coordinates Reference Manual}
\author{Author names to add}
\date{\today}

\makeindex

\begin{document}
\rowcolors[]{2}{lightgray!40}{yellow!25} %% color table rows with alternating colors
\maketitle
\tableofcontents
\clearpage

\section{Test Model - Xparser 0.17.1 Circle (x,y)-coordinates FLAME Implementation}

\subsection{Agent}

\begin{longtable}[H!]{ll}
\caption{{\bfseries List of memory variables for Agent agent.}}
\label{Table: Agent Memory}\\
\toprule 
\bfseries Name & \bfseries Description \\ \hline 
\midrule
\endfirsthead
\multicolumn{2}{c}%
{{\bfseries \tablename\ \thetable{} -- continued from previous page}} \\
\toprule
\bfseries Name & \bfseries Description \\ \hline 
\midrule
\endhead
\multicolumn{2}{r}{{\emph{Continued on next page}}} \\
\endfoot
\bottomrule
\endlastfoot
\midrule
\parbox{5cm}{\url{int} \url{id}} \index{\url{id}} & \parbox{10cm}{Agent ID.} \\
\midrule
\parbox{5cm}{\url{double} \url{x}} \index{\url{x}} & \parbox{10cm}{Agent x coordinate.} \\
\midrule
\parbox{5cm}{\url{double} \url{y}} \index{\url{y}} & \parbox{10cm}{Agent y coordinate.} \\
\midrule
\parbox{5cm}{\url{double} \url{a}} \index{\url{a}} & \parbox{10cm}{Agent x-axis offset.} \\
\midrule
\parbox{5cm}{\url{double} \url{b}} \index{\url{b}} & \parbox{10cm}{Agent y-axis offset.} \\
\end{longtable}
\begin{longtable}[H!]{ll}
\caption{{\bfseries List of functions for Agent agent.}}
\label{Table: Agent Functions}\\
\toprule 
\bfseries Name & \bfseries Description \\ \hline 
\midrule
\endfirsthead
\multicolumn{2}{c}%
{{\bfseries \tablename\ \thetable{} -- continued from previous page}} \\
\toprule
\bfseries Name & \bfseries Description \\ \hline 
\midrule
\endhead
\multicolumn{2}{r}{{\emph{Continued on next page}}} \\
\endfoot
\bottomrule
\endlastfoot
\midrule
\parbox{5cm}{\url{set_xy}} \index{\url{set_xy}} & \parbox{10cm}{} \\
\end{longtable}
\subsection{Messages}
See Table \ref{Table: Messages}.\begin{longtable}[H!]{ll}
\caption{{\bfseries List of messages.}}
\label{Table: Messages}\\
\toprule 
\bfseries Name & \bfseries Description \\ \hline 
\midrule
\endfirsthead
\multicolumn{2}{c}%
{{\bfseries \tablename\ \thetable{} -- continued from previous page}} \\
\toprule
\bfseries Name & \bfseries Description \\ \hline 
\midrule
\endhead
\multicolumn{2}{r}{{\emph{Continued on next page}}} \\
\endfoot
\bottomrule
\endlastfoot
\end{longtable}
\subsection{Constants}
\begin{longtable}[H!]{ll}
\caption{{\bfseries List of constants.}}
\label{Table: constants}\\
\toprule 
\bfseries Name & \bfseries Description \\ \hline 
\midrule
\endfirsthead
\multicolumn{2}{c}%
{{\bfseries \tablename\ \thetable{} -- continued from previous page}} \\
\toprule
\bfseries Name & \bfseries Description \\ \hline 
\midrule
\endhead
\multicolumn{2}{r}{{\emph{Continued on next page}}} \\
\endfoot
\bottomrule
\endlastfoot
\url{int} \url{set_no} \index{\url{set_no}} & \parbox{10cm}{}\\
\url{int} \url{run} \index{\url{run}} & \parbox{10cm}{}\\
\url{double} \url{const_radius} \index{\url{const_radius}} & \parbox{10cm}{}\\
\url{double} \url{const_b} \index{\url{const_b}} & \parbox{10cm}{}\\
\end{longtable}
\subsection{Datatypes}
\begin{longtable}[H!]{ll}
\caption{{\bfseries List of attributes for ADTs.}}
\label{Table: datatypes}\\
\toprule 
\bfseries Name & \bfseries Description \\ \hline 
\midrule
\endfirsthead
\multicolumn{2}{c}%
{{\bfseries \tablename\ \thetable{} -- continued from previous page}} \\
\toprule
\bfseries Name & \bfseries Description \\ \hline 
\midrule
\endhead
\multicolumn{2}{r}{{\emph{Continued on next page}}} \\
\endfoot
\bottomrule
\endlastfoot
\end{longtable}

\printindex

\end{document}
